% !PRINTED 489
\weberpage{489}

From this one obtains

\[
\begin{aligned}
u^3v^3 &= 4\,(3+\sqrt{13}), \\
u^6+v^6 &= 4\,(17+5\sqrt{13}), \\
u^3-v^3 &= \sqrt{2}\,(3+\sqrt{13}),
\end{aligned}
\]
so that, if $u^3 = \sqrt{8}\,x$ is set, the quadratic equation for $x$ follows:

\weq{17}{x^2 - \frac{3+\sqrt{13}}{2}\,(x+1) = 0, \quad f(\sqrt{-39})^3 = \sqrt{8}\,x.}

We now wish to apply our results to two problems which belong to the general theory of transformation of the seventh chapter.

\begin{center}
\textbf{\S\,132. \gesperrt{The Schläfli Modular Equation for the 23rd Degree of Transformation.}}
\end{center}

Every class equation occurs as a divisor in a great number of equations of transformation, and can therefore, if it is known in another way, be used for the calculation of equations of transformation. We take as an example the 23rd degree of transformation.

According to \S\,73 there exists, if

\weq{1}{u = f(\omega), \quad v = f(23\,\omega), \quad f\left(\frac{c+\omega}{23}\right), \quad c \equiv 0 \ (\mathrm{mod}\ 48)}

has been set, an equation between

\weq{2}{\begin{aligned}
A &= \left(\frac{u}{v}\right)^{12} + \left(\frac{v}{u}\right)^{12} \\
B &= uv + \frac{2}{uv},
\end{aligned}}

and it has already been shown in \S\,74, (14) that this equation has the form:

\weq{3}{A = B^{11} + m_1 B^{10} + \cdots + m_{10} B + m_{11},}

in which the coefficients $m_1, m_2 \ldots, m_{11}$ are rational numbers. Instead of now calculating these rational numerical coefficients, as was done there, from the developments of $u$ and $v$ in powers of $q$, we seek to determine the roots of equation (3) for $u=v$ from complex multiplication.
